\sectionTitle{Experience}
\entryHeader{Example AI Company}{AI Engineer}{2024 -- Present}
\begin{resumeBullets}
  \item Built \textbf{Project Atlas}, a multi-agent assistant for a domain workflow; designed intent routing, retrieval, and response orchestration with an LLM application framework and Python services.
  \item Delivered a document-ingestion and RAG pipeline that indexes unstructured knowledge, combines semantic retrieval with structured data, and supports multi-turn user questions.
  \item Productionized model-serving workflows with asynchronous APIs, a relational database, and observability; replace this sentence with your measurable latency, quality, cost, or adoption result.
\end{resumeBullets}

\entryHeader{Example Product Studio}{Machine Learning Engineer}{2022 -- 2024}
\begin{resumeBullets}
  \item Designed and deployed conversational AI and voice-assistant experiences for customer-facing workflows, integrating LLMs, speech services, and external APIs.
  \item Implemented hybrid retrieval that combines vector search and metadata filtering to improve access to enterprise knowledge.
  \item Engineered backend services for concurrent sessions and real-time processing; state the scale or performance target when available.
\end{resumeBullets}

\entryHeader{Example Research Lab}{AI Researcher}{2021 -- 2022}
\begin{resumeBullets}
  \item Developed a lightweight vision model for on-device classification; replace with task, F1 or accuracy, throughput, model size, and compute budget.
  \item Evaluated latency-aware architectures for edge deployment and documented reproducible experiments and ablations.
  \item Optimized inference with TensorRT and TFLite for resource-constrained devices.
\end{resumeBullets}
